% ch20-request-lifecycle
\section{رحلة الطلب الكاملة: من زرّ في الواجهة إلى صفّ في القاعدة}

في الجزء الأوّل عرفنا الطبقات الأربع، ونمط \en{CQRS} مع \en{MediatR}،
وفكرة \en{ports \& adapters}، وسلاسل المعالجة \en{(pipeline behaviors)}.
لكنّنا رأينا كلّ قطعة على حِدَة. في هذا الفصل نُركّب القطع معًا ونتتبّع
\emph{طلبًا واحدًا حقيقيًّا} من أوّله لآخره: مستخدم يفتح صفحة الإشعارات، تُرسِل
الواجهة نداءً للـ \en{API}، يعبُر النداء عدّة محطّات، تُترجَم قراءة البيانات إلى
جملة \en{SQL} مُبارَمة، ثمّ يعود الردّ ليُحدِّث الشاشة. هدفنا أن تعرف — بدقّة —
مَن يفعل ماذا في كلّ محطّة، ولماذا وُضِع بهذا الترتيب.

سنمشي «واحدة واحدة»: نبدأ من المتصفّح، ندخل الخادم، نمرّ بأنبوب الـ \en{middleware}،
ندخل \en{MediatR} وسلاسله، نصل للمعالِج ثمّ \en{EF Core}، ونرى كيف يتحوّل الكيان إلى
\en{DTO} فـ \en{JSON}. وفي الطريق نُجيب على سؤال يشغل كثيرين: لماذا يوجد أحيانًا
\en{Service} في طبقة الـ \en{Infrastructure} بدل وضع كلّ شيء داخل المعالِج؟ ونختم
بتمهيد للفرق بين مسار \en{HTTP} العاديّ ومسار الزمن الحقيقيّ \en{(SignalR)}.

% ====================================================================
\subsection{نظرة طائر: خريطة الرحلة كاملةً}

قبل التفاصيل، لِنَرَ الصورة الكبيرة. أيّ طلب قراءة يمرّ بسلسلة محطّات ثابتة، ثمّ
يرجع الردّ بنفس المسار عكسيًّا. الرسم التالي يلخّص الرحلة: الأسهم المتّصلة هي رحلة
\emph{الذهاب} (الطلب)، والسهم المتقطّع هو رحلة \emph{العودة} (الردّ).

\begin{center}\begin{LTR}
\begin{tikzpicture}[node distance=0.55cm and 0.7cm]
  \node[actor,text width=1.7cm] (br) {Browser\\(React)};
  \node[sbox,right=of br,text width=1.9cm] (mw) {Middleware\\pipeline};
  \node[pbox,right=of mw,text width=1.9cm] (ctl) {Controller};
  \node[gbox,right=of ctl,text width=1.7cm] (med) {MediatR};

  \node[tbox,below=1.15cm of med,text width=1.9cm] (beh) {Behaviors};
  \node[pbox,left=of beh,text width=1.7cm] (hnd) {Handler};
  \node[sbox,left=of hnd,text width=1.7cm] (ef) {EF Core};
  \node[dbbox,left=of ef,text width=1.7cm] (sql) {SQL\\Server};

  \draw[arr] (br) -- (mw);
  \draw[arr] (mw) -- (ctl);
  \draw[arr] (ctl) -- (med);
  \draw[arr] (med) -- (beh);
  \draw[arr] (beh) -- (hnd);
  \draw[arr] (hnd) -- (ef);
  \draw[arr] (ef) -- (sql);

  \draw[arrd] (sql.south) to[out=-90,in=-90]
    node[lbl,below] {response: entity $\rightarrow$ DTO $\rightarrow$ JSON} (br.south);
\end{tikzpicture}
\end{LTR}\end{center}

كلّ صندوق محطّة سنشرحها بالترتيب. الفكرة المحوريّة أنّ هذا المسار \emph{واحد ثابت}
لكلّ الطلبات تقريبًا؛ ما يتغيّر هو المعالِج في النهاية فقط.

\begin{spotlight}[ركّز هنا لإبهار اللجنة: مسار موحّد لكلّ طلب]
كلّ طلب في \en{ScholarPath} — قبول حجز، جلب إشعارات، نشر منشور — يمشي في
\emph{نفس} الأنبوب: \en{middleware} ثمّ \en{controller} رفيع ثمّ \en{MediatR}
ثمّ سلاسل المعالجة ثمّ معالِج واحد. إضافة ميزة جديدة تعني إضافة
\en{command/query} ومعالِجًا لها فقط — \textbf{بلا لمس البنية ولا الأنبوب}.
هذا الانتظام هو ما يجعل النظام قابلًا للنموّ دون فوضى، وهو أوّل ما تُبرِزه للجنة:
«عندنا شكل واحد لكلّ عملية، ونضيف الميزات كوحدات صغيرة مستقلّة».
\end{spotlight}

% ====================================================================
\subsection{المحطّة الأولى: الواجهة \en{(React)} تبني الطلب}

الواجهة تطبيق صفحة-واحدة \en{(SPA — Single Page Application)}. معناها أنّ المتصفّح
يُحمّل تطبيق \en{React} مرّة واحدة، وبعدها لا يُعيد تحميل الصفحة كاملةً؛ بل عند الحاجة
لبيانات جديدة \emph{يُنادي} الخادم عبر الشبكة ويستقبل \en{JSON}، ثمّ يُحدِّث جزءًا من
الشاشة. إذًا الواجهة والخادم برنامجان منفصلان يتكلّمان بلغة \en{HTTP}.

مَن يُنظّم هذه النداءات في الواجهة؟ مكتبة \en{TanStack Query}. هي طبقة فوق نداءات
الشبكة تتكفّل بثلاثة أشياء: تُطلِق النداء، وتُخزّن نتيجته في ذاكرة مؤقّتة
\en{(cache)}، وتعرف متى تُعيد الجلب. الصفحة تقول «أريد الإشعارات» فتُعطيها المكتبة
البيانات من الـ \en{cache} فورًا إن وُجدت، أو تُطلق نداءً جديدًا.

النداء الفعليّ يمرّ عبر عميل \en{axios} مشترك اسمه \code{apiClient}. أهمّ ما فيه شيءٌ
اسمه \en{interceptor}: قطعة تلتفّ حول كلّ نداء صادر وتتدخّل قبله. عندنا
\en{interceptor} على الطلب يُلصِق رأس \en{Authorization} تلقائيًّا بكلّ نداء، فلا
تحتاج كلّ صفحة أن تتذكّر ذلك.

هنشوف إيه: الـ \en{request interceptor} الحقيقيّ في \code{client.ts} — يقرأ الرمز من
مخزن الحالة ويُلصِقه كرأس \en{Bearer}.

\begin{english}
\begin{lstlisting}[language=Java]
apiClient.interceptors.request.use((config) => {
  const tokens = useAuthStore.getState().tokens;
  if (tokens?.accessToken) {
    config.headers["Authorization"] = `Bearer ${tokens.accessToken}`;
  }
  // send the active UI language so the server localises the response
  config.headers["Accept-Language"] = i18n.language || "ar";
  return config;
});
\end{lstlisting}
\end{english}

شرح الأسطر المهمّة:
\begin{itemize}
\item \code{interceptors.request.use(...)}: سجِّل دالة تُشغَّل قبل \emph{كلّ} نداء صادر.
\item نقرأ \code{tokens.accessToken} من مخزن الحالة \en{(auth store)} — رمز الدخول الحاليّ.
\item نُلصِقه في رأس \code{Authorization} بصيغة \code{Bearer <token>}؛ هذا ما يُثبِت
  للخادم هُويّة المستخدم.
\item نُرسِل أيضًا \code{Accept-Language} لتأتي الردود بلغة الواجهة (عربي أو إنجليزي).
\end{itemize}

\begin{keyp}[نقطة مهمة: الرمز يُلصَق مرّة واحدة في مكان واحد]
لولا الـ \en{interceptor} لاضطرّت كلّ صفحة أن تكتب رأس \en{Authorization} يدويًّا،
وهو تكرار خطير: نسيان واحد يعني نداءً بلا هُويّة. المركزة في مكان واحد تجعل الأمن
افتراضيًّا لا اختياريًّا.
\end{keyp}

% ────────────────────────────────────────────────────────────────
\subsubsection{ماذا لو انتهت صلاحيّة الرمز؟ الردّ \en{401} والتجديد}

رمز الدخول قصير العمر. حين ينتهي، يردّ الخادم بحالة \code{401 Unauthorized}. هنا يعمل
\en{interceptor} ثانٍ — على \emph{الردّ} هذه المرّة: إن رأى \code{401} يُحاول تجديد
الرمز عبر \code{/api/auth/refresh}، ثمّ \emph{يُعيد} النداء الأصليّ كأنّ شيئًا لم يكن.
المستخدم لا يشعر بشيء.

هنشوف إيه: قلب الـ \en{response interceptor} — التقاط \code{401}، تجديد الرمز مرّة
واحدة، ثمّ إعادة تشغيل الطلب الأصليّ.

\begin{english}
\begin{lstlisting}[language=Java]
apiClient.interceptors.response.use(
  (response) => response,
  async (error) => {
    const status = error.response?.status;
    const original = error.config;
    if (status === 401 && original && !original._retry) {
      original._retry = true;                       // guard: try refresh once
      const pending = refreshing ?? (refreshing = refreshTokens());
      const newToken = await pending;               // share ONE in-flight refresh
      if (newToken) {
        original.headers["Authorization"] = `Bearer ${newToken}`;
        return apiClient(original);                 // replay the original request
      }
      useAuthStore.getState().clear();              // refresh failed -> log out
    }
    throw new ApiError(/* ... */);
  },
);
\end{lstlisting}
\end{english}

الفكرة سطرًا بسطر:
\begin{itemize}
\item الدالة الأولى تمرّر أيّ ردّ ناجح كما هو؛ الثانية تُعالج الأخطاء.
\item \code{status === 401} مع علامة \code{\_retry}: نتأكّد أن نُجرِّب التجديد
  \emph{مرّة واحدة} فقط لا حلقة لا نهائيّة.
\item \code{refreshing ?? (...)}: لو وصلت عدّة نداءات \code{401} معًا، نتشارك عمليّة
  تجديد \emph{واحدة} بدل إطلاق عدّة عمليّات تتصادم.
\item عند نجاح التجديد نُحدِّث الرأس ونُعيد \code{apiClient(original)} — أي نُعيد
  إرسال النداء الأصليّ بالرمز الجديد.
\item عند الفشل \code{clear()} — تسجيل خروج نظيف.
\end{itemize}

\begin{note}
هذا التجديد الصامت هو سبب أنّ المستخدم يبقى داخل النظام ساعاتٍ دون أن يُطلَب منه
تسجيل الدخول ثانيةً، رغم أنّ رمز الدخول نفسه ينتهي كلّ بضع دقائق. تفاصيل الهُويّة
والرموز شرحناها في فصل التوثيق؛ هنا يهمّنا فقط كيف تتصرّف الواجهة تلقائيًّا.
\end{note}

% ====================================================================
\subsection{المحطّة الثانية: الدخول للخادم والتوجيه \en{(Routing)}}

وصل النداء الآن إلى الخادم. أوّل قرار: \emph{أيّ دالة} في أيّ \en{controller} تُعالج
هذا المسار؟ هذا هو \en{Routing} — مطابقة عنوان \en{URL} والفعل \en{(GET/POST/...)}
مع الدالة الصحيحة. في \en{ASP.NET Core} تتمّ المطابقة عبر سمات \en{(attributes)}
مكتوبة فوق الصنف والدوال.

هنشوف إيه: \code{NotificationController} — لاحظ السمات التي تُعرّف المسار، وأنّ الدالة
سطرٌ واحد يستدعي \code{mediator.Send}.

\begin{english}
\begin{lstlisting}[language={[Sharp]C}]
[ApiController]
[Route("api/notifications")]
[Authorize]
public sealed class NotificationController(IMediator mediator) : ControllerBase
{
    [HttpGet]
    public async Task<ActionResult<NotificationsPageDto>> GetMine(
        [FromQuery] int page = 1,
        [FromQuery] int pageSize = 20,
        CancellationToken ct = default)
        => Ok(await mediator.Send(new GetNotificationsQuery(page, pageSize), ct));
}
\end{lstlisting}
\end{english}

كيف يقرأ الإطار هذه السمات؟
\begin{itemize}
\item \code{[ApiController]}: يقول «هذا \en{controller} لـ \en{API}»، ويُفعّل سلوكيّات
  مريحة (ربط تلقائيّ للبيانات، ردود أخطاء موحّدة).
\item \code{[Route("api/notifications")]}: البادئة المشتركة لكلّ مسارات هذا الصنف.
\item \code{[HttpGet]} فوق \code{GetMine}: مع البادئة تعني «الفعل \en{GET} على
  \code{/api/notifications} يُنفِّذ هذه الدالة».
\item \code{[Authorize]}: لا يدخل أحد دون رمز صالح — وهنا يظهر أثر رأس \en{Bearer}
  الذي ألصقته الواجهة.
\item \code{[FromQuery]}: خُذ \code{page} و\code{pageSize} من نصّ الاستعلام في العنوان.
\end{itemize}

المهمّ أنّ الدالة \emph{رفيعة}: لا منطق عمل بداخلها. تبني كائن \code{GetNotificationsQuery}
وتُسلِّمه لـ \en{MediatR} عبر \code{mediator.Send}، وتلفّ النتيجة في \code{Ok(...)}
(أي حالة \code{200}). كلّ «كيف نجلب الإشعارات» يعيش في مكان آخر.

\begin{why}
لماذا \en{controller} رفيع لهذه الدرجة؟ لأنّ فصل «استقبال الطلب» عن «تنفيذ العمل»
يجعل المنطق قابلًا للاختبار وحده دون تشغيل خادم \en{HTTP}، ويسمح بأنبوب موحّد يلفّ
كلّ معالِج تلقائيًّا. الـ \en{controller} مجرّد بوّابة ترجمة: \en{HTTP} داخلًا، رسالة
\en{MediatR} خارجًا.
\end{why}

% ====================================================================
\subsection{المحطّة الثالثة: أنبوب الـ \en{Middleware}}

قبل أن تصل الدالة أصلًا، يمرّ الطلب عبر سلسلة من الطبقات اسمها \en{middleware}.
الـ \en{middleware} قطعة تلتفّ حول كلّ طلب: تتدخّل قبل أن يصل للـ \en{controller}،
وقد تتدخّل بعد أن يعود الردّ. تخيّلها طبقات بصلة: الطلب يدخل من الخارج للداخل حتى
المعالِج، ثمّ يخرج الردّ من الداخل للخارج عبر نفس الطبقات عكسيًّا.

\begin{center}\begin{LTR}
\begin{tikzpicture}[node distance=0.24cm]
  \node[wbox] (m1) {UseSerilogRequestLogging  —  time \& log every request};
  \node[wbox,below=of m1] (m2) {SecurityHeadersMiddleware  —  add safe HTTP headers};
  \node[wbox,below=of m2] (m3) {ExceptionHandlerMiddleware  —  errors into ProblemDetails};
  \node[wbox,below=of m3] (m4) {MaintenanceMiddleware  —  short-circuit in maintenance mode};
  \node[wbox,below=of m4] (m5) {UseAuthentication / UseAuthorization  —  who? allowed?};
  \node[pbox,below=0.34cm of m5,text width=6.6cm] (core)
    {Controller  $\rightarrow$  MediatR  $\rightarrow$  Handler};

  \draw[arr] (m1) -- (m2);
  \draw[arr] (m2) -- (m3);
  \draw[arr] (m3) -- (m4);
  \draw[arr] (m4) -- (m5);
  \draw[arr] (m5) -- (core);
  \draw[arrd] (core.east) to[out=0,in=0]
    node[lbl,right] {response travels back up} (m1.east);
\end{tikzpicture}
\end{LTR}\end{center}

الترتيب ليس عشوائيًّا. مثلًا \code{ExceptionHandlerMiddleware} موضوع \emph{مبكّرًا} حتى
يلتقط أيّ خطأ يحدث في أيّ طبقة تحته ويحوّله لردّ موحّد نظيف
\en{(ProblemDetails)}. والصيانة \code{MaintenanceMiddleware} تُقصِّر المسار إن كان
النظام في وضع صيانة فلا يُتعِب القاعدة أصلًا. والتوثيق \code{UseAuthentication} يأتي
قبل التخويل \code{UseAuthorization} لأنّه لا معنى للسؤال «هل يُسمح له؟» قبل معرفة
«مَن هو؟». سنُفرِد الفصل التالي لتفصيل هذا الأنبوب والتحقّق؛ يكفينا الآن أنّه الطبقة
التي تلفّ كلّ شيء.

\begin{idea}[الفكرة: الترتيب هو السلوك]
في أنبوب الـ \en{middleware} \emph{ترتيب} التسجيل هو نفسه ترتيب التنفيذ. تحريك سطر
واحد لأعلى أو لأسفل يُغيّر السلوك: لو وضعت معالج الأخطاء \emph{بعد} التوثيق لَما
التقط أخطاء التوثيق. القاعدة: ضَع ما يجب أن يلفّ الجميع في الخارج (مبكّرًا).
\end{idea}

% ====================================================================
\subsection{المحطّة الرابعة: \en{MediatR} وسلاسل المعالجة \en{(Behaviors)}}

سلّم الـ \en{controller} الرسالة لـ \en{MediatR}. مهمّتها إيجاد المعالِج المناسب
\en{(handler)} وتشغيله. لكن قبل أن تصل للمعالِج، تمرّ الرسالة بسلسلة أخرى داخليّة اسمها
\en{pipeline behaviors}. وهذه بالضبط فكرة \en{interceptor} لكن على مستوى الرسائل:
كلّ \en{behavior} يلتفّ حول المعالِج، يتدخّل قبله وبعده.

\begin{keyp}[نقطة مهمة: الـ \en{behavior} هو \en{interceptor} للرسائل]
عندنا نوعان من الالتفاف في النظام: \en{middleware} يلفّ طلب \en{HTTP}،
و\en{behavior} يلفّ رسالة \en{MediatR}. نفس النمط، مستوًى مختلف. الـ \en{behaviors}
عندنا تفعل أربعة أشياء عرضيّة: \textbf{التسجيل} (كم استغرق الطلب)،
\textbf{التحقّق} \en{(FluentValidation)}، \textbf{الأداء} (إنذار لو تجاوز حدًّا)،
و\textbf{التدقيق} \en{(audit)} للأوامر المُزخرَفة بـ \code{[Auditable]}.
\end{keyp}

الجميل أنّ المعالِج \emph{لا يعرف} بوجود هذه السلاسل إطلاقًا. هو يكتب منطق العمل فقط،
والتحقّق والتسجيل والتدقيق تُضاف حوله تلقائيًّا. لو أردت التحقّق من صحّة
\code{GetNotificationsQuery} تكتب صنف \en{validator} مستقلًّا، وسلسلة التحقّق تلتقطه
دون تعديل المعالِج. هذا الفصل بين «منطق العمل» و«الاهتمامات العرضيّة» هو جوهر قوّة
النمط.

% ====================================================================
\subsection{المحطّة الخامسة: المعالِج ثمّ \en{EF Core}}

وصلنا للقلب. المعالِج \en{(handler)} هو المكان الوحيد الذي يعرف \emph{كيف} تُنفَّذ هذه
العمليّة تحديدًا. لكنّه لا يكتب \en{SQL} بيده؛ يكتب استعلامًا بلغة \en{C\#} باستخدام
\en{LINQ}، و\en{EF Core} يترجمه إلى \en{SQL}. لفهم ذلك نحتاج مفهومين.

\begin{idea}[الفكرة: \en{IQueryable} والتنفيذ المؤجَّل]
\code{IQueryable} ليس بيانات؛ هو \emph{وصف} لاستعلام لم يُنفَّذ بعد. حين تكتب
\code{.Where(...)} ثمّ \code{.OrderBy(...)} أنت تبني \emph{شجرة تعبير} تصف ما تريد،
لا أكثر. لا يذهب شيء للقاعدة حتى تطلب النتيجة فعليًّا. هذا اسمه \en{deferred
execution} — التنفيذ المؤجَّل. اللحظة الحاسمة هي نداء مثل \code{ToListAsync} أو
\code{FirstOrDefaultAsync} أو \code{CountAsync}: عندها \emph{فقط} تُترجَم الشجرة
لـ \en{SQL} وتُرسَل وتُنفَّذ.
\end{idea}

هنشوف إيه: معالِج \code{GetNotificationsQueryHandler} الحقيقيّ. راقب أين يُبنى الاستعلام
وأين يُنفَّذ فعلًا.

\begin{english}
\begin{lstlisting}[language={[Sharp]C}]
public async Task<NotificationsPageDto> Handle(
    GetNotificationsQuery request, CancellationToken ct)
{
    var userId = currentUser.UserId
        ?? throw new ForbiddenAccessException("Not authenticated.");

    // IQueryable: nothing has run yet, just a description of the query
    var baseQuery = db.Notifications.Where(
        n => n.RecipientUserId == userId && n.Channel == NotificationChannel.InApp);

    var total = await baseQuery.CountAsync(ct);       // runs: SELECT COUNT(*)

    var items = await baseQuery
        .OrderByDescending(n => n.CreatedAt)
        .Skip((request.Page - 1) * pageSize)
        .Take(pageSize)
        .Select(n => new NotificationDto(n.Id, n.TitleAr, n.IsRead, n.CreatedAt))
        .ToListAsync(ct);                             // <-- SQL runs HERE

    return new NotificationsPageDto(items, request.Page, pageSize, total, unread);
}
\end{lstlisting}
\end{english}

النقاط المفتاحيّة:
\begin{itemize}
\item \code{db.Notifications.Where(...)} يُنتِج \code{baseQuery} من نوع
  \code{IQueryable} — \emph{لم يُنفَّذ بعد}، ولا استعلام ذهب للقاعدة.
\item \code{.Where(n => n.RecipientUserId == userId ...)}: مرشّح — نُريد إشعارات هذا
  المستخدم فقط، وقناة العرض داخل التطبيق فقط.
\item \code{.OrderByDescending(...).Skip(...).Take(...)}: ترتيب وتقسيم صفحات — كلّها
  ما زالت \emph{وصفًا} يُضاف للشجرة.
\item \code{.Select(n => new NotificationDto(...))}: إسقاط \en{(projection)} —
  نطلب من القاعدة الأعمدة التي نحتاجها فقط لا الصفّ كاملًا.
\item \code{await ...ToListAsync(ct)}: \emph{هنا} تُترجَم الشجرة كلّها لجملة \en{SELECT}
  واحدة، وتُرسَل، وتعود النتائج. سطر \code{CountAsync} قبله نفّذ \code{SELECT COUNT(*)}
  منفصلًا.
\end{itemize}

كيف تبدو الترجمة؟ \en{EF Core} يمتلك «مزوِّدًا» \en{(provider)} لـ \en{SQL Server}
يمشي في شجرة التعبير ويبني الجملة المقابلة. والأهمّ: أيّ قيمة متغيّرة (مثل
\code{userId}) لا تُلصَق في نصّ الجملة، بل تُمرَّر كـ \emph{بارامتر} منفصل.

\begin{center}\begin{LTR}
\begin{tikzpicture}[node distance=0.7cm]
  \node[srcbox,text width=5.3cm,align=left,minimum height=1.6cm] (cs)
    {\textbf{C\# / LINQ}\\[2pt]
     db.Notifications\\
     \ \ .Where(n => n.RecipientUserId == userId)\\
     \ \ .OrderByDescending(n => n.CreatedAt)};
  \node[dbbox,right=1.4cm of cs,text width=5.3cm,align=left,minimum height=1.6cm] (sql)
    {\textbf{SQL (parameterized)}\\[2pt]
     SELECT Id, TitleAr, IsRead, CreatedAt\\
     FROM Notifications\\
     WHERE RecipientUserId = @p0\\
     ORDER BY CreatedAt DESC};
  \draw[arr] (cs) -- node[lbl,above,align=center]
    {EF Core provider\\translates} (sql);
  \node[dbox,below=0.5cm of sql,text width=5.0cm]
    {@p0 = userId  (a bound parameter, never string-concatenated)};
\end{tikzpicture}
\end{LTR}\end{center}

القيمة \code{userId} تُرسَل في خانة البارامتر \code{@p0} (اسمها الفعليّ عند
\en{EF Core} أطول قليلًا، لكن الفكرة واحدة). النصّ لا يحوي القيمة أبدًا. لهذا أثرٌ أمنيّ
كبير جدًّا.

\begin{spotlight}[ركّز هنا لإبهار اللجنة: الترجمة الآمنة \en{LINQ} إلى \en{SQL}]
لأنّ \en{EF Core} يُمرِّر كلّ القيم كبارامترات منفصلة، فإنّ \textbf{حقن \en{SQL}
\en{(SQL Injection)} يصبح مستحيلًا افتراضيًّا} — دون أن يكتب المطوّر سطرًا واحدًا
للحماية. لو أدخل مستخدمٌ خبيثٌ نصًّا مثل \code{'; DROP TABLE ...} كقيمة، فإنّه يُعامَل
\emph{بيانةً} في خانة البارامتر لا \emph{أمرًا}، فلا يُنفَّذ. هذه نقطة قويّة أمام
اللجنة: «أماننا ضدّ الحقن ليس حِيلة أضفناها، بل خاصيّة بنيويّة نحصل عليها مجّانًا لأنّنا
لا نبني \en{SQL} نصيًّا أبدًا».
\end{spotlight}

\begin{secnote}
القاعدة الذهبيّة: لا تبنِ استعلامًا بدمج نصوص \en{(string concatenation)}. طالما بقيت
داخل \en{LINQ}، فالبَرمَتة تلقائيّة. الخطر الوحيد يظهر لو لجأ أحدٌ لـ \en{SQL} خام
مكتوب يدويًّا بدمج مدخلات المستخدم — وهذا ممنوع في مشروعنا.
\end{secnote}

% ====================================================================
\subsection{المحطّة السادسة: الرجوع — كيان إلى \en{DTO} إلى \en{JSON}}

عادت النتائج من القاعدة على شكل كائنات في الذاكرة. لكنّنا لا نُرسِل الكيان
\en{(entity)} الداخليّ للواجهة مباشرةً؛ نُرسِل \en{DTO} — كائن نقل بيانات يحمل فقط ما
تحتاجه الشاشة. في مثالنا رأينا الإسقاط \code{Select(n => new NotificationDto(...))}
يبني الـ \en{DTO} مباشرةً داخل الاستعلام. في حالات أخرى نستخدم \en{AutoMapper} ليُحوِّل
الكيان إلى \en{DTO} بقاعدة تحويل مُعرَّفة مسبقًا، فلا نكتب النسخ حقلًا حقلًا.

بعدها يتولّى الإطار تسلسل الـ \en{DTO} إلى نصّ \en{JSON} ويضعه في جسم الردّ مع حالة
\code{200}. يعود هذا الردّ صعودًا عبر أنبوب الـ \en{middleware} (مسجِّلًا زمنه في
الطريق) حتى يخرج للشبكة.

\begin{why}
لماذا \en{DTO} وليس الكيان مباشرةً؟ لثلاثة أسباب: لا نُسرِّب حقولًا داخليّة أو حسّاسة
للواجهة؛ ولا نربط شكل الـ \en{API} بشكل جداول القاعدة (نستطيع تغيير أحدهما دون كسر
الآخر)؛ ونُرسِل بيانات أقلّ على الشبكة لأنّ الـ \en{DTO} يحمل ما يلزم فقط.
\end{why}

في الواجهة تستقبل \en{TanStack Query} الـ \en{JSON}، تحوّله لكائنات، وتُخزّنه في
الـ \en{cache} تحت مفتاح الاستعلام. أيّ مكوّن يطلب نفس البيانات يحصل عليها فورًا من
الذاكرة. وعند حدوث تغيير (مثل تعليم إشعار مقروءًا) نُبطِل \en{(invalidate)} المفتاح،
فتُعيد المكتبة الجلب تلقائيًّا وتُحدِّث الشاشة. بهذا تُغلَق الدائرة: من زرٍّ في الواجهة،
إلى صفٍّ في القاعدة، ورجوعًا إلى بكسل على الشاشة.

% ====================================================================
\subsection{سؤالك الصريح: لماذا يوجد \en{Service} في الـ \en{Infrastructure} ولا نضع كلّ شيء في المعالِج؟}

سؤال ممتاز ويستحقّ إجابةً صريحة. رأينا أنّ المعالِج يقرأ من القاعدة عبر \en{EF Core}
مباشرةً، فلماذا لا نفعل \emph{كلّ} شيء داخله؟ الجواب: قراءة القاعدة عمليّة نظيفة
ومملوكة لنا، أمّا بعض الأعمال فتحتاج التعامل مع \emph{عالم خارجيّ} لا نتحكّم فيه —
وهذه نعزلها في \en{Service} داخل الـ \en{Infrastructure} (وهو تنفيذ لـ \en{port}
مُعرَّف في طبقة \en{Application}). أربعة أسباب:

\begin{center}
\begin{tabular}{L{4.2cm} L{8.6cm}}
\toprule
\textbf{\color{spPrimary}السبب} & \textbf{\color{spPrimary}الشرح} \\
\midrule
اتّصال \en{I/O} خارجيّ &
  الدفع عبر \en{Stripe}، إرسال البريد، رفع ملفّ إلى \en{Blob}، غرفة \en{ACS} — كلّها
  نداءات شبكة لطرف ثالث، بمنطق إعادة محاولة ومفاتيح سرّيّة. لا يليق حشوها في معالِج. \\
إعادة الاستخدام &
  \code{StripeService} أو \code{IEmailService} يُستدعى من معالِجات كثيرة. نكتب المنطق
  مرّة ونحقنه حيث نشاء بدل تكراره. \\
تبديل المزوّد بالإعداد &
  المعالِج يعتمد على \en{port} (واجهة) لا على تنفيذ بعينه. نستبدل التنفيذ من الإعداد:
  \code{StubMeetingService} محليًّا و\code{AzureCommunicationMeetingService} إنتاجيًّا،
  دون لمس المعالِج. \\
قابليّة الاختبار &
  في الاختبار نحقن \en{stub} مزيّفًا للـ \en{port} فنختبر منطق المعالِج دون لمس
  \en{Stripe} أو الشبكة إطلاقًا — سريع وحتميّ. \\
\bottomrule
\end{tabular}
\end{center}

القاعدة العمليّة: \textbf{القراءة/الكتابة البسيطة للقاعدة تبقى في المعالِج عبر
\en{EF Core}؛ وأيّ تعامل مع نظام خارجيّ يُعزَل خلف \en{port} ويُنفَّذ كـ \en{Service}
في الـ \en{Infrastructure}}. هكذا يبقى المعالِج مركّزًا على «قصّة العمل»، وتبقى
تفاصيل الشبكة والمزوّدين قابلةً للتبديل والاختبار.

\begin{qa}[سؤال متوقّع: أليس هذا طبقات زائدة بلا داعٍ؟]
لا، لأنّ العزل يشتري لنا شيئين ملموسين: نُشغّل النظام محليًّا بلا حساب \en{Stripe}
أو \en{Azure} (عبر تنفيذات \en{stub})، ونكتب اختبارات لا تلمس الشبكة فتكون سريعة
وموثوقة. الكلفة سطر واجهة واحد؛ العائد استقلالٌ عن مزوّدين نتحكّم فيهم بالإعداد لا
بالكود.
\end{qa}

% ====================================================================
\subsection{تمهيد: مسار \en{HTTP} مقابل مسار الزمن الحقيقيّ \en{(SignalR)}}

كلّ ما سبق كان مسار \en{HTTP}: العميل \emph{يسأل} والخادم \emph{يجيب}، ثمّ يُغلَق
الاتّصال. هذا مثاليّ لطلبٍ صريح مثل «أعطني الإشعارات». لكن ماذا لو أردنا العكس: أن
\emph{يدفع} الخادم خبرًا للعميل دون أن يسأل — رسالة دردشة وصلت الآن، أو إشعار جديد
ظهر؟ في \en{HTTP} وحده لا يستطيع الخادم بدء الكلام.

هنا يأتي \en{SignalR}: قناة \emph{مفتوحة ثنائيّة الاتّجاه} (عبر \en{WebSockets}) تبقى
حيّةً بين العميل والخادم، فيستطيع أيّ طرف الإرسال في أيّ لحظة. الفرق الجوهريّ:

\begin{center}
\begin{tabular}{L{4.0cm} L{4.3cm} L{4.3cm}}
\toprule
 & \textbf{\color{spPrimary}\en{HTTP} (طلب/ردّ)} & \textbf{\color{spPrimary}\en{SignalR} (زمن حقيقيّ)} \\
\midrule
مَن يبدأ الكلام؟ & العميل دائمًا & أيّ طرف \\
عمر الاتّصال & لحظيّ (يُغلَق بعد الردّ) & مفتوح ودائم \\
الاتّجاه & باتّجاه واحد لكلّ نداء & ثنائيّ \\
مناسب لِـ & جلب/حفظ بيانات & إشعارات ودردشة مباشرة \\
\bottomrule
\end{tabular}
\end{center}

في \code{Program.cs} نربط الاثنين جنبًا إلى جنب: \code{app.MapControllers()} لمسارات
\en{HTTP}، و\code{app.MapHub(...)} لقنوات \en{SignalR} مثل \code{/hubs/notifications}
و\code{/hubs/chat}. نفس التطبيق يخدم المسارين. تفاصيل هذا المسار الحيّ — كيف يُبنى
الاتّصال، وكيف يُبثّ الإشعار للمستخدم المعنيّ فقط — موضوع فصلٍ قادمٍ مستقلّ؛ هنا فقط
غرسنا البذرة: طلبٌ يسأل، وقناةٌ تدفع.

% ====================================================================
\begin{summary}[الخلاصة في سطور]
\begin{itemize}
\item الواجهة \en{(React + TanStack Query)} تبني النداء، و\en{interceptor} على العميل
  يُلصِق \en{Authorization: Bearer} تلقائيًّا، ويجدّد الرمز عند \code{401} ويُعيد الطلب.
\item التوجيه \en{(routing)} عبر \code{[ApiController]} و\code{[HttpGet]} يوصّل المسار
  للدالة الصحيحة؛ والـ \en{controller} رفيع: سطرٌ واحد \code{mediator.Send}.
\item الـ \en{middleware} سلسلة تلفّ كلّ طلب (تسجيل، أمان، أخطاء، صيانة، توثيق)؛
  ترتيبها هو سلوكها.
\item الـ \en{MediatR pipeline behaviors} هي \en{interceptor} للرسائل: تحقّق وتسجيل
  وأداء وتدقيق حول المعالِج دون أن يعلم.
\item المعالِج يكتب \en{LINQ} على \code{IQueryable} (تنفيذ مؤجَّل)؛ لا شيء يُنفَّذ حتى
  \code{ToListAsync}؛ و\en{EF Core} يترجم لـ \en{SQL} مُبارَم يمنع حقن \en{SQL}.
\item الردّ: كيان إلى \en{DTO} إلى \en{JSON}، والواجهة تُحدِّث \en{cache}‎ها؛ فتُغلَق
  الدائرة من الزرّ إلى الصفّ ورجوعًا.
\item ما يتعامل مع عالم خارجيّ (دفع/بريد/تخزين) يُعزَل في \en{Service} خلف \en{port}
  للإعادة والتبديل والاختبار.
\item \en{HTTP} للطلب/الردّ، و\en{SignalR} للدفع الحيّ ثنائيّ الاتّجاه — وهو فصل قادم.
\end{itemize}
\end{summary}
